\documentclass[10pt,a4paper]{amsart}
\usepackage[margin=19mm]{geometry}
\usepackage{amsmath,amssymb,mathtools}
\usepackage[scr=rsfs,cal=euler]{mathalfa}
\usepackage{xfrac}
\usepackage[unicode,hidelinks,bookmarks=false]{hyperref}
\newcommand{\ie}{i.e.\ }
\newcommand{\cf}{cf.\ }
\newcommand{\resp}{respectively\ }
\newcommand{\Subsection}{\subsection}
\providecommand{\nobreakdash}{\nobreak\hbox{-}}
\setlength{\emergencystretch}{2em}
\multlinegap0pt
\usepackage{bm}
\usepackage{amssymb}
\usepackage{float}
\usepackage{xcolor}
\newtheorem{lemma}{Lemma}
\newcommand{\CC}{\mathbb{C}}
\newcommand{\RR}{\mathbb{R}}
\title{Quantum Fisher information matrix via its classical counterpart from random measurements}
\author{Jianfeng Lu, Kecen Sha}
\date{Source: arXiv:2509.08196v4 (CC BY 4.0)}
\begin{document}
\maketitle

\noindent\emph{Source and licence.} Quantum Fisher information matrix via its classical counterpart from random measurements. Version arXiv:2509.08196v4 (CC BY 4.0), distributed by arXiv under the Creative Commons Attribution 4.0 International licence, http://creativecommons.org/licenses/by/4.0/, as recorded in the arXiv licence metadata for this exact version and shown on the abstract page https://arxiv.org/abs/2509.08196v4. Full licence notice: \texttt{licenses/arxiv-2509.08196-CC-BY-4.0.txt}.\par\smallskip

\noindent\textit{Editorial scope.} Counted unit: the article's lemma lem:keyLip (source0.tex lines 827--831). The article's complete original proof of it is delivered with it (source0.tex lines 832--896, one \texttt{proof} environment of 65 lines). No mathematics is written, repaired or re-derived: the statement and the proof are the article's own lines, in the article's own notation, and no retained line is substituted or re-flowed.  No further source-local context is needed: every cross-reference of every retained line resolves inside the two delivered spans.  Nothing was omitted from any retained span, so the package carries no omission record.  No retained line contains a citation, so the wrapper carries no bibliography and no bibliography entry.  No retained line contains a figure environment, an image-inclusion command or drawing code, so no image is delivered and the package declares no asset dependency.  Editorial formatting only, in addition: an AMS \texttt{amsart} preamble that restates the article's own declarations and macros used by the retained lines, namely the environments \texttt{lemma} on separate counters, together with the standard packages those lines need.  The retained environments print in this document as Lemma 1 in order of appearance, whereas the article prints the same items with the numbering of its own full document.  The complete native source of the article as distributed in the arXiv e-print for this exact version is bundled unchanged as \texttt{sources/184803/source0.tex}.
\begin{lemma}
\label{lem:keyLip}
    Given a parameter $\bm\theta$ and a unit vector $\bm r$ that has no zero entry, for any $U,V\in \mathrm U(N)$ such that $U^*\psi_{\bm\theta}=V^*\psi_{\bm\theta}=\bm r$, we have 
    \[\|P_{S^U}(\bm z_{\bm\theta})-P_{S^V}(\bm z_{\bm\theta})\|_2\leq \frac{\pi}{2}\|U-V\|_F,\quad \|P_{S^U}(\bm z_{\bm\theta})-P_{S^V}(\bm z_{\bm\theta})\|_F\leq \frac{\sqrt 2\pi}{2}\|U-V\|_F.\]
\end{lemma}

\begin{proof}
    Denote $\bm r=(r_1,\cdots,r_N)^\top$ then $\|\Phi(U) D_k\Phi(U)^\top J\bm{z}_{\bm\theta}\|_2^2=|r_k|^2\neq 0$. Let $\bm{w}_k=\Phi(iU\tilde D_k\bm r)$,then
\begin{equation}
    P(\Phi(U) D_k\Phi(U)^\top J\bm{z}_{\bm\theta})=P(\bm{w}_k)=|r_k|^{-2}\bm{w}_k\bm{w}_k^\top.
\end{equation}
 Let $\delta U,\delta\bm{w}_k$ be the small variation of $U,\bm{w}_k$ respectively. By their definitions, it is easy to check that $\delta\bm{w}_k=\Phi(ir_k\delta U\bm e_k)$. Moreover, the tangent space of $\mathrm U(N)$ at $U$ is $\{UY:Y^*=-Y\}$. Hence, $\delta U=U\delta Y+O(\delta^2)$ where $\delta Y=\delta\cdot Y$ is a skew-Hermitian matrix.  $\delta\bm{w}_k=\Phi(ir_kU\delta Y\bm e_k)+O(\delta^2)$, so we have 
\begin{equation}
\label{eq:Project error}
    P(\bm{w}_k+\delta\bm{w}_k)-P(\bm{w}_k)=\frac{1}{|r_k|^2}(\bm{w}_k\delta\bm{w}_k^\top+\delta\bm{w}_k\bm{w}_k^\top)+O(\delta^2).
\end{equation}

Denote $\tilde{\bm{w}}_k=\frac{\bm{w}_k}{|r_k|},\delta\tilde{\bm{w}}_k=\frac{\delta\bm{w}_k}{|r_k|}$. Then $\|\tilde{\bm{w}}_k\|_2=1,\|\delta\tilde{\bm{w}}_k\|_2=\delta\|Y\bm e_k\|_2+O(\delta^2),\delta\tilde{\bm{w}}_k^\top\tilde{\bm{w}}_k=O(\delta^2)$, 
\begin{equation}
    \frac{1}{|r_k|^2}(\bm{w}_k\delta\bm{w}_k^\top+\delta\bm{w}_k\bm{w}_k^\top)=(\tilde{\bm{w}}_k\delta\tilde{\bm{w}}_k^\top+\delta\tilde{\bm{w}}_k\tilde{\bm{w}}_k^\top)+O(\delta^2).
\end{equation}

Let $W=[\tilde{\bm{w}}_1,\cdots,\tilde{\bm{w}}_N, J\tilde{\bm{w}}_1,\cdots,J\tilde{\bm{w}}_N]$, $\delta W=[\delta\tilde{\bm{w}}_1,\cdots,\delta\tilde{\bm{w}}_N,0,\cdots,0]^\top$. Then 
\begin{equation}
\label{eq:error matrix form}
   \sum\limits_{k=1}^N(\tilde{\bm{w}}_k\delta\tilde{\bm{w}}_k^\top+\delta\tilde{\bm{w}}_k\tilde{\bm{w}}_k^\top)=W\delta W+(\delta W)^\top W^\top.
\end{equation}

 Since $\delta\bm{w}_k^\top\bm{w}_l=\Phi(ir_kU\delta Y\bm e_k)^\top\Phi(ir_lU\bm e_l)+O(\delta^2)$, and $\langle ir_kU\delta Y\bm e_k,ir_l U\bm e_l\rangle_{\CC}=r_k\bar r_l\delta Y_{lk}$, we have
 \begin{equation}
     \delta\tilde{\bm{w}}_k^\top\tilde{\bm{w}}_l=\delta\operatorname{Re}(\frac{r_k\bar r_l}{|r_kr_l|}Y_{lk})+O(\delta^2).
 \end{equation}

Similarly, we have 
\begin{equation}
    \delta\tilde{\bm{w}}_k^\top J\tilde{\bm{w}}_l=\delta\operatorname{Im}(\frac{r_k\bar r_l}{|r_kr_l|}Y_{lk})+O(\delta^2).
\end{equation}

Let $\tilde Y_{kl}=\frac{r_k\bar r_l}{|r_kr_l|}Y_{lk}$. Note that $\tilde Y$ is also skew-Hermitian, so $\operatorname{Re}(\tilde Y)$ is anti-symmetric and $\operatorname{Im}(\tilde Y)$ is symmetric. Using the above results for  $\delta\tilde{\bm{w}}_k^\top \tilde{\bm{w}}_l,\delta\tilde{\bm{w}}_k^\top J\tilde{\bm{w}}_l$, we can calculate each entry of the matrix $\delta W W+W^\top(\delta W)^\top$. We get
\begin{equation}
    \delta W W+W^\top(\delta W)^\top=\delta \begin{bmatrix}
    0&\operatorname{Im}(\tilde Y)\\
   \operatorname{Im}(\tilde Y)&0
\end{bmatrix}+O(\delta^2),
\end{equation}
Note that due to the anti-symmetry of $\operatorname{Re}(\tilde Y)$, the top-left block is zero.
\begin{align}
\label{ineq:norm bound}
   \|\delta WW+W^\top(\delta W)^\top\|_2&\leq\delta\|\operatorname{Im}(\tilde Y)\|_2+O(\delta^2)\leq \delta\|\tilde Y\|_2+O(\delta^2).
\end{align}

Now we can bound $\|P_{S^{U+\delta U}}-P_{S^U}(\bm{z}_{\bm\theta})\|_2$ as follows:
\begin{align}
\label{eq:final key norm bound}
    \|P_{S^{U+\delta U}}(\bm{z}_{\bm\theta})-P_{S^U}(\bm{z}_{\bm\theta})\|_2&=\|\sum\limits_{k=1}^{N}(P(\bm{w}_k+\delta\bm{w}_k)-P(\bm{w}_k))\|_2\notag\\
    &=\| W\delta W+(\delta W)^\top W^\top\|_2+O(\delta^2)\notag \\
    &=\|\delta W W+W^\top(\delta W)^\top\|_2+O(\delta^2)\notag\\
    &\leq \delta\|Y\|_2+O(\delta^2).
\end{align}
The first equality follows from (\ref{eq:Project error})$\sim$(\ref{eq:error matrix form}). The second equality follows from the fact that $W$ is an orthogonal matrix in $\RR^{2N\times 2N}$. The last inequality follows from (\ref{ineq:norm bound}) and $\|Y\|_2=\|\tilde Y\|_2$.

Note that $\|\delta U\|_2=\|U\delta Y\|_2+O(\delta^2)=\delta\|Y\|_2+O(\delta^2)$, so we have $\|P_{S^{U+\delta U}}(\bm{z}_{\bm\theta})-P_{S^U}(\bm{z}_{\bm\theta})\|_2\leq\|\delta U\|_2+O(\delta^2)$. Fix $U_0,U_1\in \mathrm U(N)$ such that $U_0^*\psi_{\bm\theta}=\bm e_1,U_1^*\bm r=\bm e_1$, then for every $U$ such that  $U^*\psi_{\bm\theta}=\bm r$, there exists a unique $X\in \mathrm U(N-1)$ such that $U=U_0\begin{bmatrix}
    1&0\\
    0& X
\end{bmatrix}U_1^*$, and vice versa. Denote this map as $U=g(X)$, then it is easy to see that $\|g(X_1)-g(X_2)\|_2=\|X_1-X_2\|_2$ and $\|g(X_1)-g(X_2)\|_F=\|X_1-X_2\|_F$. Let $U=g(X)$ and $V=g(Z)$.
For every $X,Z$ on the compact manifold $\mathrm U(N-1)$, the geodesic distance $d(X,Z)$ between $X,Z$ is no more than $\frac{\pi}{2}\|X-Z\|_F$. Take $M+1$ matrices on the geodesic between $X,Z$, namely $X_0,X_1,\cdots,X_{M}$ such that $X_0=X,X_{M}=Z$ and $\|X_{i+1}-X_i\|_F=O(M^{-1})$, then we have  
\[\|P_{S^U}(\bm z_{\bm\theta})-P_{S^V}(\bm z_{\bm\theta})\|_2\leq\sum\limits_{k=0}^{M-1}\|P_{S^{g(X_k)}}(\bm z_{\bm\theta})-P_{S^{g(X_{k+1})}}(\bm z_{\bm\theta})\|_2\leq \sum\limits_{k=0}^{M-1}\|X_{k+1}-X_{k}\|_F+M\cdot O(M^{-2}).\]
Let $M\to\infty$, we have $\|P_{S^U}(\bm z_{\bm\theta})-P_{S^V}(\bm z_{\bm\theta})\|_2\leq d(X,Z)\leq \frac{\pi}{2} \|X-Z\|_F=\frac{\pi}{2}\|U-V\|_F$.

In a similar way to (\ref{eq:final key norm bound}), we have $\|P_{S^U}(\bm z_{\bm\theta})-P_{S^{U+\delta U}}(\bm z_{\bm\theta})\|_F\leq \sqrt{2}\delta \|Y\|_F+O(\delta^2)$. By the same geodesic argument, we know that $\|P_{S^U}(\bm z_{\bm\theta})-P_{S^V}(\bm z_{\bm\theta})\|_F\leq \sqrt{2}d(X,Z)\leq\frac{\sqrt{2}\pi}{2}\|U-V\|_F$. 
\end{proof}

\end{document}
